\documentclass[journal]{IEEEtran}
\usepackage{amsmath,amssymb}
\usepackage{graphicx}
\usepackage{booktabs}
\usepackage{multirow}
\usepackage{cite}
\usepackage[colorlinks=false,pdfborder={0 0 0}]{hyperref}

\begin{document}

\title{From Prediction to Decision: Uncertainty-Aware Deferral for Reliable
Medical Image Segmentation}

\author{Author Names Withheld}

\maketitle

\begin{abstract}
Deep learning segmentation models are increasingly deployed in clinical
workflows, yet most pipelines lack systematic policies for deciding which
predictions to trust. We study this as a selective prediction problem,
comparing how two uncertainty estimation methods (Monte Carlo Dropout and
Test-Time Augmentation) and three deferral strategies (global thresholding,
image-adaptive thresholding, confidence-aware deferral) interact in a decision
pipeline for retinal vessel segmentation. Evaluated on DRIVE with cross-dataset
testing on STARE and CHASE\_DB1, we find that: (1) TTA produces uncertainty with
AUROC 0.881 for error prediction versus 0.768 for MC Dropout, at 6.3$\times$
lower latency; (2) confidence-aware deferral, which jointly weighs uncertainty
and prediction confidence, reduces segmentation error by 55\% while deferring
only 12\% of pixels, compared to 26\% error reduction at 33\% deferral for
global thresholding; (3) post-hoc temperature scaling does not improve deferral
performance despite marginally reducing calibration error; and (4) uncertainty
quality, measured by AUROC, transfers more robustly than segmentation accuracy
under domain shift. These findings support evaluating uncertainty methods by
their decision-level consequences rather than calibration metrics.
\end{abstract}

\begin{IEEEkeywords}
Uncertainty estimation, selective prediction, medical image segmentation,
deferral, retinal vessel segmentation, calibration
\end{IEEEkeywords}

\section{Introduction}

Medical image segmentation models are evaluated by aggregate metrics: Dice
coefficient, IoU, AUC. These averages hide per-image and per-pixel variability.
A model with 0.78 Dice may produce near-perfect segmentations on some images
and clinically unacceptable outputs on others. For deployment, what matters is
not the average but the policy: which predictions should the system trust, and
which should it flag for human review?

This question connects two research threads. Uncertainty estimation methods,
from Bayesian approximations~\cite{gal2016dropout} to ensembles~\cite{lakshminarayanan2017simple},
produce maps that correlate with model error. Selective prediction, studied
primarily in classification~\cite{geifman2017selective,el2010foundations},
formalizes the idea that a model should abstain when confidence is low. In
segmentation, the connection between these ideas remains underdeveloped. Most
work generates uncertainty maps and computes calibration metrics without asking
how those maps should drive downstream decisions.

We address this gap. We do not propose new architectures or uncertainty
estimators. Instead, we study the decision layer that sits on top of existing
methods: given an uncertainty map and a prediction, what deferral policy should
the system follow?

\subsection{Contributions}

\begin{enumerate}
\item A systematic evaluation of MC Dropout and TTA as inputs to deferral
policies, showing that TTA produces substantially better uncertainty for
decision-making (Unc-AUROC 0.881 vs.\ 0.768) while being 6.3$\times$ faster.

\item A confidence-aware deferral strategy that jointly uses uncertainty
magnitude and prediction confidence, achieving 4.6$\times$ better deferral
efficiency (error reduction per deferred pixel) than global thresholding.

\item Evidence that temperature scaling does not translate to better deferral
performance, questioning the assumption that calibration improves clinical
utility.

\item Cross-dataset evaluation showing that uncertainty quality transfers more
robustly than segmentation accuracy, with Unc-AUROC remaining above 0.83 on
STARE and CHASE\_DB1 despite a 5\% Dice drop.

\item A risk-coverage framework for evaluating uncertainty methods by their
ability to improve segmentation through selective abstention.
\end{enumerate}

\section{Related work}

\subsection{Uncertainty estimation}

MC Dropout~\cite{gal2016dropout} approximates Bayesian inference by retaining
dropout at test time. Multiple stochastic passes yield a distribution over
predictions; the variance or mutual information quantifies epistemic
uncertainty. MC Dropout is widely used in medical
imaging~\cite{nair2020exploring,jungo2018effect,roy2019bayesian} because it
requires no architectural changes.

Test-Time Augmentation~\cite{wang2019aleatoric,ayhan2018test} applies
deterministic transforms to the input and measures prediction disagreement.
Mehrtash et al.~\cite{mehrtash2020confidence} compared MC Dropout and TTA for
brain tumors and found TTA uncertainty correlated more strongly with error. We
confirm this for retinal vessels and extend it to the deferral setting.

Deep ensembles~\cite{lakshminarayanan2017simple} train multiple models with
different initializations. They produce well-calibrated
predictions~\cite{ovadia2019can} but multiply training and inference cost.

\subsection{Calibration}

Temperature scaling~\cite{guo2017calibration} learns a single scalar to rescale
logits, optimizing ECE~\cite{naeini2015obtaining}. The assumption is that
calibrated probabilities support better decisions. Laves et
al.~\cite{laves2020well} observed that calibration and selective prediction can
trade off against each other in classification. We extend this finding to dense
pixel-level deferral.

\subsection{Selective prediction and deferral}

Selective prediction has a long history~\cite{chow1970optimum,el2010foundations}.
Geifman and El-Yaniv~\cite{geifman2017selective} formalized risk-coverage
curves for deep classifiers. In medical imaging, Jungo et
al.~\cite{jungo2020analyzing} studied uncertainty thresholds for brain lesion
segmentation. Madras et al.~\cite{madras2018predict} and Mozannar and
Sontag~\cite{mozannar2020consistent} formalized learning-to-defer with known
expert models. Our approach requires no expert model and operates on any
existing segmentation pipeline without retraining.

\section{Methods}

\subsection{Problem formulation}

Let $\mathbf{x} \in \mathbb{R}^{H \times W \times C}$ be an input image,
$\mathbf{y} \in \{0,1\}^{H \times W}$ the ground truth, and $f_\theta$ a
segmentation model producing $\hat{\mathbf{p}} \in [0,1]^{H \times W}$. An
uncertainty method produces $\mathbf{u} \in \mathbb{R}_{\geq 0}^{H \times W}$.
A deferral policy $\delta$ maps $(\mathbf{u}, \hat{\mathbf{p}})$ to binary
accept/defer decisions. We evaluate by the error reduction ratio:
$\text{ERR} = (e_{\text{before}} - e_{\text{after}}) / e_{\text{before}}$.

\subsection{MC Dropout}

With $T$ stochastic passes, the mean prediction is $\bar{p}_{ij} =
\frac{1}{T}\sum_t \hat{p}_{ij}^{(t)}$. Uncertainty is mutual information:
\begin{equation}
\text{MI}_{ij} = H(\bar{p}_{ij}) - \frac{1}{T}\sum_t H(\hat{p}_{ij}^{(t)})
\end{equation}
where $H(p) = -p\log p - (1-p)\log(1-p)$.

\subsection{Test-Time Augmentation}

With $K$ geometric transforms $\{\tau_k\}$, the mean is $\bar{p}_{ij} =
\frac{1}{K}\sum_k [\tau_k^{-1} \circ f_\theta \circ \tau_k](\mathbf{x})_{ij}$.
Uncertainty is prediction variance across augmented outputs.

\subsection{Deferral policies}

\textbf{Global:} $\delta_{ij} = \mathbb{1}[u_{ij} \leq \tau]$, threshold
$\tau$ fit on validation set.

\textbf{Adaptive:} $\delta_{ij} = \mathbb{1}[u_{ij} \leq Q_\alpha(\mathbf{u}_n)]$,
per-image percentile threshold.

\textbf{Confidence-aware:} Define $c_{ij} = 2|\hat{p}_{ij} - 0.5|$ and
$s_{ij} = u_{ij} \cdot (1 - c_{ij})$. Defer if $s_{ij} > \tau_s$. Pixels
that are uncertain \emph{and} near the decision boundary are deferred first.

\subsection{Temperature scaling}

$\hat{p}^{\text{cal}} = \sigma(\text{logit}(\hat{p}) / T)$ with $T$
optimized by minimizing BCE on validation logits.

\section{Experimental setup}

U-Net with ResNet-34 encoder, pretrained on ImageNet. Training: DRIVE (20
images), 256$\times$256 patches with vessel-aware sampling, 80 epochs, Adam
($\text{lr} = 10^{-4}$), Dice+BCE loss (positive weight 4.0). MC Dropout:
$T = 30$, $p = 0.3$. TTA: $K = 6$ (identity, flips, 90/180/270$^\circ$
rotations). Cross-dataset: STARE (20 images), CHASE\_DB1 (28 images).

\section{Results}

\subsection{Uncertainty method comparison}

TTA: Dice 0.768, Unc-AUROC 0.881, AUCC (Dice) 0.846, 0.13\,s/image.
MC Dropout: Dice 0.764, Unc-AUROC 0.768, AUCC 0.801, 0.82\,s/image.
Ensemble ($N = 5$): Dice 0.771, Unc-AUROC 0.833, 3.90\,s/image.

TTA produces better uncertainty at lower cost. The Dice difference (0.768
vs.\ 0.764) is not significant ($p = 0.142$); the Unc-AUROC difference is
($p < 0.001$).

\subsection{Deferral strategy comparison}

\begin{table}[t]
\centering
\caption{Deferral results on DRIVE.}
\small
\begin{tabular}{llcccc}
\toprule
Uncertainty & Deferral & ERR & Def.\ \% & ERR/Def.\ \% \\
\midrule
MC Dropout & Global & 26.4\% & 33.4\% & 0.79 \\
MC Dropout & Conf-aware & 48.9\% & 10.9\% & 4.49 \\
TTA & Global & 42.2\% & 20.0\% & 2.11 \\
TTA & Adaptive & 79.5\% & 25.0\% & 3.18 \\
TTA & Conf-aware & 55.2\% & 12.1\% & 4.56 \\
\bottomrule
\end{tabular}
\end{table}

Confidence-aware deferral achieves 4.5$\times$ better efficiency (ERR per
deferred percentage point) than global thresholding.

\subsection{Calibration analysis}

Temperature scaling: $T = 1.19$ (MC), $T = 1.35$ (TTA). ECE changes:
0.0382 $\to$ 0.0377 (MC), 0.0366 $\to$ 0.0393 (TTA). Unc-AUROC changes:
0.768 $\to$ 0.749 (MC, degradation), 0.882 $\to$ 0.881 (TTA, no change).
Calibration does not improve and can degrade the uncertainty-error correlation
that deferral depends on.

\subsection{Cross-dataset evaluation}

DRIVE $\to$ STARE: Dice 0.764 $\to$ 0.727 ($-$5\%), Unc-AUROC 0.768 $\to$
0.846 ($+$10\%). DRIVE $\to$ CHASE\_DB1: Dice 0.726, Unc-AUROC 0.835.
Uncertainty quality \emph{improves} on out-of-distribution data because errors
concentrate in regions the model also finds uncertain.

\subsection{Ablation studies}

MC Dropout passes ($T \in \{5, 10, 20, 30\}$): Unc-AUROC saturates at $T = 10$
(0.836 vs.\ 0.838 at $T = 30$). Runtime scales linearly (0.24\,s to 1.42\,s).

Cross-validation (5-fold): Dice $0.780 \pm 0.004$, Unc-AUROC
$0.768 \pm 0.057$.

\section{Discussion}

TTA outperforms MC Dropout because geometric augmentations directly probe
prediction fragility at boundaries and thin structures, where errors
concentrate. MC Dropout's weight-space perturbations produce more diffuse
uncertainty that does not localize to error-prone pixels.

Confidence-aware deferral succeeds by recognizing that uncertainty alone is
not sufficient: a pixel can be uncertain yet correctly predicted. The score
$s = u \cdot (1-c)$ concentrates deferrals on the highest-risk pixels,
those that are both uncertain and near the decision boundary.

Temperature scaling fails to help because it optimizes a global distributional
property (ECE) while deferral requires a local discriminative property
(uncertainty-error separation). A model optimized for calibration distributes
errors across confidence bins; a model useful for deferral concentrates errors
at low confidence.

\section{Conclusion}

Uncertainty should be evaluated by its downstream decision consequences, not
calibration metrics. TTA with confidence-aware deferral provides the best
combination of uncertainty quality, deferral efficiency, and computational cost.
Future work should extend to multi-class segmentation and validate with
human-in-the-loop studies.

\bibliographystyle{IEEEtran}
\bibliography{../references}

\end{document}
